\begin{qfdefinition}[Commodities as investments: their features, the forms of access, and their risk, return and diversification characteristics]\label{def:commodity-features}
\qffield{Statement} \emph{Features.} Commodity futures differ from stocks and bonds in three ways: ``(1) commodity futures are derivative securities; they are not claims on long-lived corporations; (2) they are short maturity claims on real assets; (3) unlike financial assets, many commodities have pronounced seasonality in price levels and volatilities'' \citep[p.~2]{s16-r34}. Every commodity trades on a spot market, and a commodity trade may be physical (delivery) or financial (cash settlement), unlike trades in stocks and bonds \citep[pp.~1--2]{physical-geman-commodities-2005}. Some commodities are storable and some, such as electricity, are not (Gorton and Rouwenhorst, p.~2; Geman, p.~251).

\emph{Forms of access.} Financial investors rarely hold commodities physically, because of storage costs, and take exposure through futures; even precious-metal investors mostly use futures or funds \citep[p.~365]{s7-ang}. The other routes are index funds and ETFs, swaps on commodity indices (\hyperref[meth:commodity-swaps-index-construction]{swaps and indices}), commodity producers' shares, managed-futures funds (Ang, p.~365, on commodity trading advisors) and, for timber and farmland, the land itself (\hyperref[meth:land-timber-farmland]{land, timber and farmland}).

\emph{Investment characteristics.}
\begin{itemize}
\item \emph{Return}: a collateralised futures position earns the spot return, the roll return and the interest on its collateral (\hyperref[meth:commodity-futures-return-components]{return components}); commodities pay no income, so a long-run return above cash must come from a risk premium for insuring hedgers or from rebalancing (Gorton and Rouwenhorst, p.~2); \citep[pp.~10--12, 24--25]{wide-erb-harvey-commodity}.
\item \emph{Risk}: high volatility --- 20.7 percent a year for the GSCI from 1986 to 2011, with a raw Sharpe ratio of 0.43 against 0.66 for stocks (Ang, pp.~366--367). Volatility rises when inventories are low and so, unlike equities, tends to rise with prices (Geman, p.~28).
\item \emph{Diversification}: the GSCI's correlations with stocks and bonds were 14 and $-9$ percent over 1986--2011 (Ang, p.~367); Gorton and Rouwenhorst find collateralised futures negatively correlated with stocks and bonds from 1959 to 2004, partly because they behave differently over the business cycle (abstract).
\item \emph{Inflation}: only energy has hedged inflation, with a GSCI correlation of 29 percent driven by energy (23--26 percent); agriculture's was $-4$ percent and gold's 1 percent (Ang, pp.~366--367).
\end{itemize}

\qffield{Practice} \emph{Spot holder against futures holder} (Gorton and Rouwenhorst's example, p.~6): oil trades at \$30, the one-month future at \$25 and the expected spot price is \$27. If oil ends at \$27, the holder of physical oil has lost \$3 while the futures buyer has gained \$2, the risk premium; if it ends at \$26, the futures buyer gains \$1, the premium less the \$1 surprise. A futures investor is paid for bearing price risk, not for the commodity's appreciation.

\qffield{Conditions} Figures are historical and index-specific; the GSCI is about 80 percent energy (Ang, p.~367). Gorton and Rouwenhorst and Erb and Harvey are cited by the printed pages of their NBER working papers (w10595, revised March 2006; w11222, March 2005, titled ``The Tactical and Strategic Value of Commodity Futures''); the 2006 \emph{Financial Analysts Journal} versions were not read. Gorton and Rouwenhorst's abstract gives a sample to March 2004 while their figures run to December 2004.

\qffield{Pitfalls} Expecting futures returns to follow spot prices. Treating commodities as a general inflation hedge.

\qffield{See also} \hyperref[meth:commodity-sectors-participants]{sectors}; \hyperref[def:commodity-indices-risk-return]{index record}; \hyperref[meth:alternatives-allocation-approaches]{allocating to alternatives}.

\end{qfdefinition}

\begin{qfmethod}[Commodity sectors: their characteristics, life cycles from production to consumption, valuation against equities and bonds, and the participants in futures markets]\label{meth:commodity-sectors-participants}
\qffield{Statement} \emph{Sectors and life cycles} \citep[pp.~143--144, 169, 201, 251--252]{physical-geman-commodities-2005}:
\begin{itemize}
\item \emph{Agriculture} (grains, oilseeds, softs, livestock): supply is carry-in stocks plus the year's production plus imports; demand is domestic use and exports; carry-over links crop years. The life cycle is planting, growing and harvest, so supply is set once a year and depends on weather; the stocks-to-use ratio (typically 20--40 percent) signals tightness, and USDA planting and acreage reports in March and June move prices (pp.~143--144). Livestock adds a feeding cycle of several months; produce is perishable.
\item \emph{Metals}: mined, concentrated and refined; the traded ingot is ``approximately halfway through the life cycle'' that runs from ore underground to semi-fabricated products and alloys (p.~169). Supply responds slowly because mines take years to develop; metals store easily, and gold is held largely as a store of value.
\item \emph{Energy}: crude oil is produced upstream, refined into products and consumed; the market has become a financial market with horizons beyond ten years and ``marker'' grades such as WTI underpin a much larger derivatives market (p.~201). Natural gas is costly to store and seasonal.
\item \emph{Electricity} is ``mostly non-storable'', so supply and demand must balance in real time and prices spike from \$30 to \$1,000 a megawatt-hour or more (pp.~251--252).
\end{itemize}

\emph{Valuation against equities and bonds.} Stocks and bonds are claims on cash flows, valued as ``the discounted value of cash flows over very long horizons''; commodity futures ``do not raise resources for firms'' but let firms insure the value of their outputs or inputs, and their investors are paid ``for bearing the risk of short-term commodity price fluctuations'' \citep[p.~2]{s16-r34}. A commodity's spot price is set by current supply, demand and inventory, and its futures price by the spot price, storage costs, interest, the convenience yield of holding inventory and any risk premium \citep[pp.~24--25, 33--34]{physical-geman-commodities-2005}; \citep[pp.~365, 370]{s7-ang}. There is no discounted-cash-flow value to compare with the price.

\emph{Participants} \citep[pp.~6--10]{physical-geman-commodities-2005}: \emph{hedgers}, producers and consumers locking in prices (an airline buying fuel futures), who still bear basis risk; \emph{speculators}, including investors and hedge funds, who take price exposure and provide most liquidity with hedgers; \emph{arbitrageurs} trading spot against futures; and exchanges, clearing houses and futures commission merchants. The US Commodity Futures Trading Commission regulates the markets (p.~9).

\qffield{Practice} \emph{Reading a grain balance sheet} (Geman's figures, p.~144): ending wheat stocks of 540 million bushels against use of 2.7 billion give a stocks-to-use ratio of 20 percent, at the tight end of the usual range; a weather shock with steady demand would then raise prices sharply, and the spot--futures spread would move toward backwardation (\hyperref[def:contango-backwardation-stocks]{contango and backwardation}).

\qffield{Conditions} Geman's book describes markets as of 2004--2005; exchange names have since changed through mergers.

\qffield{Pitfalls} Applying equity valuation multiples to commodities. Ignoring seasonality and storability when comparing sectors.

\qffield{See also} \hyperref[def:commodity-features]{features}; \hyperref[meth:commodity-futures-return-components]{return components}.

\end{qfmethod}
